
\documentclass[a4paper,11pt]{article}
\usepackage{fontspec}
\setmainfont{Noto Sans}
\usepackage[french]{babel}
\usepackage{microtype}
\usepackage[top=9mm,bottom=8mm,left=11mm,right=11mm]{geometry}
\usepackage{xcolor}
\usepackage{tikz}
\usetikzlibrary{arrows.meta,calc,positioning}
\usepackage[most]{tcolorbox}
\pagestyle{empty}

\definecolor{Orange}{HTML}{D8842F}
\definecolor{OrangeDark}{HTML}{9A5716}
\definecolor{OrangeLight}{HTML}{FFF1DF}
\definecolor{Border}{HTML}{E7C7A6}
\definecolor{Ink}{HTML}{172A32}
\definecolor{Grey}{HTML}{64757D}
\definecolor{Red}{HTML}{D52F35}
\definecolor{RedLight}{HTML}{FCEBED}

\tcbset{
 enhanced,boxrule=.6pt,arc=2.8mm,
 left=3.7mm,right=3.7mm,top=1.75mm,bottom=1.75mm,
 boxsep=1mm,coltext=Ink
}
\newtcolorbox{orangebox}[1][]{colback=OrangeLight,colframe=Border,#1}
\newtcolorbox{warnbox}[1][]{colback=RedLight,colframe=Red,
 left=2.7mm,right=2.7mm,top=1.25mm,bottom=1.25mm,#1}

\begin{document}
\begin{tikzpicture}[remember picture,overlay]
\draw[draw=Border,line width=.9pt,rounded corners=1.7mm]
 ([xshift=1mm,yshift=-1mm]current page.north west)
 -- ([xshift=-1mm,yshift=-1mm]current page.north east)
 -- ([xshift=-1mm,yshift=1mm]current page.south east)
 -- ([xshift=1mm,yshift=1mm]current page.south west)--cycle;
\fill[Orange] ([xshift=1mm,yshift=-1mm]current page.north west)
 rectangle ([xshift=5.2mm,yshift=1mm]current page.south west);
\end{tikzpicture}

\vspace*{1mm}\hspace*{6mm}
{\color{OrangeDark}\fontsize{25}{27}\selectfont\bfseries FUTUR DE L'INDICATIU}

\vspace{0.5mm}\hspace*{6mm}
{\color{Grey}\large Le futur de l'indicatif}

\vspace{2.5mm}\hspace*{6mm}
\begin{minipage}{0.91\textwidth}

\begin{orangebox}
\textbf{\color{OrangeDark}CONSTRUCCION}\\[1mm]
\centering
{\Large \textbf{INFINITIU + TERMINASONS DEL FUTUR}}\\[1mm]
\small
Exemple : \textbf{CANTAR}\\[1mm]
\textbf{cantarai}\quad \textbf{cantaràs}\quad \textbf{cantarà}\quad
\textbf{cantarem}\quad \textbf{cantaretz}\quad \textbf{cantaran}
\end{orangebox}

\vspace{2.5mm}
\begin{center}
\begin{tikzpicture}[>=Latex]
\draw[Ink!65,line width=.8pt] (0,0)--(15.1,0);
\node[below=2mm] at (2.0,0) {\small PASSAT};
\node[below=2mm] at (7.5,0) {\small ARA};
\node[below=2mm] at (13.5,0) {\small FUTUR};
\draw[Orange,line width=2.5pt,->] (8.0,0)--(13.0,0);
\draw[OrangeDark,line width=1.1pt,->] (7.5,1.0)--(7.5,.18);
\node[above=3.8mm] at (7.5,1.0) {\bfseries\color{OrangeDark}ARA};
\node[above=3mm,align=center,text width=6.2cm] at (10.5,0)
{\small\bfseries accion futura, projècte o prediccion};
\end{tikzpicture}
\end{center}

\vspace{1.5mm}
\begin{orangebox}
\textbf{\color{OrangeDark}A QUÉ SERVÍS ?}\\[1mm]
\begin{itemize}\setlength{\itemsep}{0.4mm}
\item parlar d'una accion situada \textbf{après lo moment present} ;
\item exprimir un \textbf{projècte}, una decision o una promessa ;
\item far una \textbf{prediccion} ;
\item presentar un fach futur considerat coma segur o previst.
\end{itemize}
\end{orangebox}

\vspace{2mm}
\begin{orangebox}
\textbf{\color{OrangeDark}EXEMPLES}\\[1mm]
\textbf{Demà, partirai d'ora.}\\
\textbf{Aquel estiu, trabalharem ensemble.}\\
\textbf{Quand vendrà, li parlarai.}\\
\textbf{Lo temps cambiarà lèu.}
\end{orangebox}

\vspace{2mm}
\begin{minipage}{0.485\textwidth}
\begin{orangebox}[height=3.25cm]
\textbf{\color{OrangeDark}TERMINASONS}\\[1mm]
\centering
\textbf{-ai \quad -às \quad -à}\\
\textbf{-em \quad -etz \quad -an}\\[1mm]
\raggedright
L'infinitiu servís de basa per fòrça vèrbs regulars.
\end{orangebox}
\end{minipage}\hfill
\begin{minipage}{0.485\textwidth}
\begin{warnbox}[height=3.25cm]
\textbf{\color{Red}ATENCION}\\[1mm]
Certans vèrbs an un \textbf{radical irregular} al futur.\\[1mm]
Per exemple :\\
\textbf{èsser → serai}\\
\textbf{aver → aurai}\\
\textbf{far → farai}
\end{warnbox}
\end{minipage}

\vspace{2mm}
\begin{orangebox}
\textbf{\color{OrangeDark}A RETÉNER}\\[1mm]
Lo futur de l'indicatiu permet de situar una accion \textbf{après ara}. Es frequent per parlar
d'un projècte, d'una prediccion, d'una promessa o d'un eveniment futur.
\end{orangebox}

\vspace{1.5mm}\centering
{\small\color{Grey}\textbf{Indicis frequents :} demà, lèu, lèu-lèu, la setmana que ven, l'an que ven, mai tard, un jorn, quand...}

\vspace{1mm}
{\scriptsize\color{Grey}Referéncia : occitan languedocian, graphia classica, vèrb'Òc / Lo Congrès permanent de la lenga occitana.}

\end{minipage}
\end{document}
